% concatenated from main_v2.tex
\documentclass[12pt]{article}

\usepackage[margin=1.05in]{geometry}
\usepackage{amsmath,amssymb,amsthm}
\usepackage{mathrsfs}
\usepackage{graphicx}
\usepackage{url}
\usepackage{hyperref}


\newtheorem{theorem}{Theorem}[section]
\newtheorem{proposition}[theorem]{Proposition}
\newtheorem{lemma}[theorem]{Lemma}
\newtheorem{corollary}[theorem]{Corollary}
\newtheorem{conjecture}[theorem]{Conjecture}

\theoremstyle{definition}
\newtheorem{definition}[theorem]{Definition}

\theoremstyle{remark}
\newtheorem{remark}[theorem]{Remark}

\newcommand{\R}{\mathbb R}
\newcommand{\Z}{\mathbb Z}
\newcommand{\E}{\mathbb E}
\newcommand{\Var}{\operatorname{Var}}
\newcommand{\Cov}{\operatorname{Cov}}
\newcommand{\Tr}{\operatorname{Tr}}
\newcommand{\Id}{\operatorname{Id}}
\newcommand{\HS}{\mathrm{HS}}
\newcommand{\op}{\mathrm{op}}
\newcommand{\dd}{\,\mathrm d}
\newcommand{\ip}[2]{\left\langle #1,#2\right\rangle}
\newcommand{\norm}[1]{\left\lVert #1\right\rVert}
\newcommand{\abs}[1]{\left\lvert #1\right\rvert}
\newcommand{\cL}{\mathcal L}
\newcommand{\cG}{\mathcal G}
\newcommand{\cK}{\mathcal K}
\newcommand{\cE}{\mathcal E}
\newcommand{\cS}{\mathcal S}
\newcommand{\tauM}{\tau_M}
\newcommand{\tauL}{\tau_L}
\newcommand{\sym}{\operatorname{sym}}

\hypersetup{
  pdftitle={Optimal MM* bounds for convex bodies},
  pdfauthor={Pierre Bizeul},
}

\title{Optimal MM* bounds for convex bodies}

\author{Pierre Bizeul}

\date{}

\begin{document}

\maketitle

\begin{abstract}
Let $K\subset\R^n$ be a convex body in isotropic position. We prove the
optimal mean-width estimate
\[
  M^*(K)\leq C\sqrt{n\log n}.
\]
The main new ingredient is a geometric inequality relating the Gaussian
mean of the support function to its mean under the uniform measure on
$K$, obtained through a heat-flow argument. Combined with the
Gaussian-log-concave comparison of Eldan and Lehec and the newly available
dimension-free bound on the third-moment parameter $\kappa_n$, this
yields the result. The boundedness of $\kappa_n$ also makes the mean-norm
estimate of Bizeul and Klartag sharp. Combining both estimates yields
\[
  M(K)M^*(K)\leq C\log n,
\]
extending Pisier's $MM^*$ estimate to non-symmetric convex bodies and to the isotropic position.
\end{abstract}



\section{Introduction}

Throughout the paper, we work in $\R^n$ with $n\geq2$. A convex body
$K\subset\R^n$ is a compact convex subset with non-empty interior. We
shall denote by $X_K$ the random vector uniformly distributed in $K$.
More generally, we say that a random vector $X$ is log-concave if it
has a density of the form $e^{-V}$, with
$V:\R^n\longrightarrow\R\cup\{+\infty\}$ a convex function. We say
that a random vector $X$ is isotropic if it is centered and its
covariance matrix is the identity. A convex body $K$ will be called
isotropic if the random vector $X_K$ is. Namely,
\[
  \int_K x\dd x=0,
  \qquad
  \frac1{|K|}\int_K x\otimes x\dd x=\Id,
\]
where $|\cdot|$ denotes Lebesgue volume. Note that this convention,
which we retain throughout the paper, is sometimes called
probabilistic isotropy and is slightly different from volume
isotropy, which requires the volume of $K$ to be one and its covariance
matrix to be scalar. Any convex body can be made isotropic after
applying a suitable invertible affine transformation. In the
terminology of convex geometry, an invertible affine image of a convex
body is sometimes called a ``position''.

The support function of $K$ is defined by
\[
  h_K(x)=\sup_{y\in K}\langle x,y\rangle,
  \qquad x\in\R^n.
\]
If $K$ contains the origin in its interior, we define its gauge by
\[
  \norm{x}_K
  =\inf\left\{t>0:\ \frac{x}{t}\in K\right\}.
\]
Under the same assumption, the polar body of $K$ is defined as the
body whose gauge is $h_K$, namely
\[
  K^\circ
  =
  \left\{
    x\in\R^n:\
    \sup_{y\in K}\langle x,y\rangle\leq1
  \right\}.
\]
Note that $\norm{\cdot}_K$ and $h_K$ define genuine norms only when
$K$ is origin-symmetric.

The mean-norm and mean-width are two fundamental geometric parameters
associated with a convex body $K$, defined respectively by
\[
  M(K)=\E\norm{\Theta}_K,
  \qquad
  M^*(K)=\E h_K(\Theta),
\]
where $\Theta$ is distributed uniformly on the sphere
$\mathbb S^{n-1}$. These two quantities are dual to one another since,
tautologically,
\[
  M(K^\circ)=M^*(K).
\]

It is not too difficult to provide lower bounds for $M$ and $M^*$.
Indeed, for a fixed volume, Urysohn's inequality states that the
mean-width is minimized by the Euclidean ball:
\begin{equation}\label{eq_157}
  M^*(K)
  \geq
  \left(\frac{|K|}{|B_2^n|}\right)^{1/n} := \operatorname{vrad}(K).
\end{equation}
Here and below, $B_p^n$ denotes the unit ball of $\ell_p^n$.
Furthermore, when $K$ contains the origin in its interior, expressing
its volume in polar coordinates and using Jensen's inequality gives
the corresponding estimate
\begin{equation}\label{eq_166}
  M(K)
  \geq
  \left(\frac{|B_2^n|}{|K|}\right)^{1/n}.
\end{equation}
In particular, the scale-invariant quantity
\[
  \ell(K):=M(K)M^*(K)
\]
satisfies
\[
  \ell(K)\geq1.
\]
The
question of upper-bounding these quantities in an appropriate position
is central to convex geometry.

The main result of this paper is the following optimal estimate on the
mean-width in isotropic position.

\begin{theorem}\label{thm_186}
Let $K$ be an isotropic convex body in $\R^n$. Then
\begin{equation}\label{eq_188}
  M^*(K)\leq C\sqrt{n\log n},
\end{equation}
where $C>0$ is a universal constant.
\end{theorem}

Theorem~\ref{thm_186} is optimal, up to the value of the
constant $C$, as follows from the example of the cross-polytope. Before
the present work, the best known estimate in isotropic position was due
to Emanuel Milman~\cite{Milman}:
\[
  M^*(K)\lesssim\sqrt n\,\log^2 n.
\]
Earlier estimates were obtained in
\cite{HartzoulakiThesis,PivovarovMeanWidth}; as observed in
\cite{BrazitikosGiannopoulosValettasVritsiou}, they may also be deduced
from Paouris's work~\cite{PaourisConcentration}.

After the first version of this paper appeared, Brazitikos and Pandis~\cite{BrazitikosPandis} obtained,
for origin-symmetric isotropic bodies, the intermediate estimates
\[
  M(K)\lesssim\frac{\log n}{\sqrt n},
  \qquad
  M^*(K)\lesssim\sqrt n\,\log n,
\]
using different methods, based on centroid bodies.

Finally, in independent and simultaneous work, Paouris and
Pathak~\cite{PaourisPathak} proved the sharp affine estimate
\[
  \inf_{T\in GL_n}
  \frac{M^*(TK)}{\operatorname{vrad}(TK)}
  \lesssim \sqrt{\log n}.
\]
However, their proof does not use the isotropic position, but rather assumes that $K^\circ$ is in Bobkov position \cite{BobkovMPosition} : namely the gaussian measure restricted to $K^\circ$ is assumed to have scalar covariance. In particular, it does not imply the existence of a position where both $M$ and $M^*$ are controlled, i.e it does not yield an $MM^*$ bound.

As explained below, both the result of Bizeul-Klartag \cite{BizeulKlartag} as well as Theorem \ref{thm_186} also implies an affine optimal bound on $M^*$ by applying the $M$-estimate
to the polar body in isotropic position and using the Bourgain–Milman reverse Santal´o inequality. On the other
hand, Theorem \ref{thm_186} gives the same affine conclusion directly from the isotropic position
once the slicing theorem is used. Both of these arguments rely
on the newly available dimension-free estimate for $\kappa_n$ (see below for a discussion on the role of this parameter). Remarkably, the Paouris–Pathak
argument does not require any such bound.


The proof of Theorem~\ref{thm_186} combines two ingredients.
The first is the following optimal comparison between Gaussian and
log-concave vectors in the gauge order.

\begin{theorem}\label{thm_206}
Let $X$ be an isotropic log-concave random vector in $\R^n$, and let
$\norm{\cdot}$ be a gauge. Then
\begin{equation}\label{eq_209}
  \frac{\E\norm{G}}{C\sqrt{\log n}}
  \leq \E\norm{X}
  \leq C\sqrt{\log n}\,\E\norm{G},
\end{equation}
where $G\sim N(0,\Id)$ is a standard Gaussian vector and $C>0$ is a
universal constant.
\end{theorem}

The left-hand side of this inequality was established by Bizeul and
Klartag~\cite{BizeulKlartag}, while the right-hand side was proved by
Eldan and Lehec~\cite{EldanLehec}. In both cases, the original
statements involved an additional parameter $\kappa_n$, related to
the third-order moment tensor of isotropic log-concave vectors.
Therefore, the statement given above depends on a dimension-free
estimate for $\kappa_n$, which is now available~\cite{Letwin}. We
postpone the discussion of this parameter and of its dimension-free
estimate until the end of the introduction.

The second ingredient is the following new inequality.

\begin{theorem}\label{thm_230}
Let $K$ be a centered convex body in $\R^n$. Then
\begin{equation}\label{eq_232}
  \bigl(\E h_K(G)\bigr)^2
  \leq
  2n\,\E h_K(X_K),
\end{equation}
where $G\sim N(0,\Id)$.
\end{theorem}

Theorem~\ref{thm_230} does not require an isotropic normalization
and admits a self-contained heat-flow proof. It is the main new
ingredient in the proof of Theorem~\ref{thm_186}. Taken
together, Theorems~\ref{thm_206} and~\ref{thm_230}
imply Theorem~\ref{thm_186}, as explained in
Section~\ref{sec_940}.

Theorem~\ref{thm_186} complements the corresponding mean-norm
estimate of Bizeul and Klartag, which, combined with the
dimension-free bound on $\kappa_n$, yields the following.

\begin{theorem}[Bizeul-Klartag~\cite{BizeulKlartag}]
\label{thm_252}
Let $K$ be an isotropic convex body in $\R^n$. Then
\begin{equation}\label{eq_254}
  M(K)
  \leq
  C\sqrt{\frac{\log n}{n}},
\end{equation}
where $C>0$ is a universal constant.
\end{theorem}

The bound is sharp, as follows from the example of the cube. Prior to \cite{BizeulKlartag}, mean-norm estimates in isotropic position were obtained in 
\cite{GiannopoulosMilman2,GiannopoulosStavrakakisTsolomitisVritsiou, VritsiouRegularEllipsoids}.

As mentionned previously, Theorem~\ref{thm_252} already yields the
sharp affine bounds for $M$ and $M^*$. Indeed, after the
normalization $|K|=|B_2^n|$, it gives an affine position in which
\[
  M(K)\lesssim\sqrt{\log n}.
\]
Applying the same estimate to a polar body placed in isotropic
position, and using the Bourgain--Milman reverse Santal\'o
inequality~\cite{BourgainMilman,ArtsteinAvidanGiannopoulosMilman}, gives
another affine position with bounded bolume ratio in which
\[
  M^*(K)\lesssim\sqrt{\log n}.
\]
Thus the optimal affine estimates for $M$ and $M^*$ already follow from
the Bizeul-Klartag estimate together with the dimension-free bound on
$\kappa_n$. The
only volume estimate required is the elementary fact that an isotropic body $K$ satisfies
$\operatorname{vrad}(K)\lesssim\sqrt n$, which is the "easy" direction of the slicing problem. The solution of the slicing problem (\cite{KlartagLehecSlicing, BizeulSlicing} and the very recent \cite{BrazitikosSlicing}) in fact implies that $\operatorname{vrad}(K)\simeq\sqrt n$. This allows to show that the $M^*$ bound in isotropic position of Theorem \ref{thm_186} yields an affine optimal bound. See the remark after the proof of
Corollary~\ref{cor_319}.

The new point of Theorem~\ref{thm_186} is the
optimal bound for $M^*$ in the isotropic position. Together with
Theorem~\ref{thm_252}, this yields an optimal isotropic $MM^*$ estimate.

\begin{corollary}\label{cor_265}
Let $K$ be an isotropic convex body in $\R^n$. Then
\begin{equation}\label{eq_267}
  \ell(K)=M(K)M^*(K)\leq C\log n.
\end{equation}
\end{corollary}

This should be compared with the landmark result of Pisier (see \cite{Pisier}), which
states that for every origin-symmetric convex body $K\subset\R^n$,
there exists an invertible linear transformation $T\in GL_n$ such
that
\begin{equation}\label{eq_276}
  \ell(TK)\leq C\log n.
\end{equation}
Corollary~\ref{cor_265} extends this estimate to
arbitrary convex bodies, without any symmetry assumption, and shows
that one may choose the position to be isotropic.
Before the present work, $MM^*$ estimates in the non-symmetric setting were pioneered by Rudelson \cite{Rudelson} with a bound of order $n^{1/3}$, up to logarithmic factors. More recently,
Bizeul and Klartag established an optimal mean-norm bound, up to the value of $\kappa_n$, and combined it with Emanuel Milman's mean-width estimate to obtain $\ell(K) \leq C\log^3 n$
in isotropic position.

We now include a
discussion of the optimality of the results presented above, beyond the isotropic
normalization.

\paragraph{Optimality of the estimates.}

The three estimates
\[
  M(K)\lesssim\sqrt{\frac{\log n}{n}},
  \qquad
  M^*(K)\lesssim\sqrt{n\log n},
  \qquad
  \ell(K)\lesssim\log n
\]
are all optimal in isotropic position. This remains true even if one
restricts attention to origin-symmetric convex bodies. As mentioned
above, the cube and the cross-polytope saturate the first two
inequalities, respectively. The third one is saturated by a product of
a cube and a cross-polytope, each factor living in half the dimension.
Details are provided in Section~\ref{sec_1045}.

For the individual quantities $M$ and $M^*$, a stronger statement holds. The estimates are in fact affine optimal : they cannot be
improved by choosing a different affine position.

\begin{corollary}\label{cor_319}
The mean-width and mean-norm estimates are affine optimal. Namely
\begin{equation}\label{eq_321}
  \sup_K\inf_{T,a} M(TK+a)
  \simeq\sqrt{\log n},
  \qquad
  \sup_K\inf_{T,a} M^*(TK+a)
  \simeq\sqrt{\log n}.
\end{equation}
Here, the suprema run over all convex bodies $K\subset\R^n$ satisfying
$|K|=|B_2^n|$, and the infima run over all $T\in SL_n$ and
$a\in\R^n$ such that $0\in\operatorname{int}(TK+a)$. The same
conclusions remain valid when the suprema are restricted to
origin-symmetric convex bodies.

\end{corollary}

The situation for the product $\ell(K)$ is subtler. Since $\ell(K)$
is scale invariant, no volume normalization is required. In the class
of all convex bodies, one has

\begin{corollary}\label{cor_340}
Corollary \ref{cor_265} is affine optimal. Namely
\begin{equation}\label{eq_342}
  \sup_K\inf_{T,a} \ell(TK+a)
  \simeq\log n,
\end{equation}
where the supremum runs over all convex bodies $K\subset\R^n$, and the
infimum runs over all $T\in GL_n$ and $a\in\R^n$ such that
$0\in\operatorname{int}(TK+a)$.
\end{corollary}

Thus, Corollary~\ref{cor_265} is optimal even if one is allowed to
choose an arbitrary affine position. In this sense, the extension of
Pisier's estimate to non-symmetric convex bodies is optimal. In contrast, if one restricts attention to origin-symmetric convex
bodies, it is believed that Pisier's estimate may be improved; see
e.g. \cite[Section~6.8]{ArtsteinAvidanGiannopoulosMilman}.

\begin{conjecture}\label{conj_358}
For every origin-symmetric convex body $K\subset\R^n$, there exists an
invertible linear transformation $T\in GL_n$ such that
\begin{equation}\label{eq_361}
  \ell(TK)\leq C\sqrt{\log n},
\end{equation}
where $C>0$ is a universal constant.
\end{conjecture}

The conjectural order is the best possible, since it is saturated by
the cube, or equivalently by the cross-polytope. Since the logarithmic
estimate in Corollary~\ref{cor_265} is optimal even within the
class of origin-symmetric convex bodies, the conjecture above cannot
hold in general if one insists on using the isotropic position.

All the optimality statements above, together with the examples
establishing them, are discussed in greater detail in
Section~\ref{sec_1045}.

\paragraph{The parameter $\kappa_n$.}
We conclude the introduction by discussing the parameter $\kappa_n$
and its newly available dimension-free bound, used in
Theorem~\ref{thm_206}. Define
\begin{equation}\label{eq_381}
  \kappa_n
  =
  \sup_X\sup_{\theta\in\mathbb S^{n-1}}
  \norm{
    \E\bigl[\langle X,\theta\rangle X\otimes X\bigr]
  }_{\HS},
\end{equation}
where the first supremum runs over all isotropic log-concave random
vectors $X$ in $\R^n$, and $\norm{\cdot}_{\HS}$ denotes the
Hilbert-Schmidt norm. A dimension-free estimate for $\kappa_n$ was
recently obtained in~\cite{Letwin}. For the reader's convenience, we
explain how a very short modification of the argument of Chen and
Klartag~\cite{ChenKlartag} recovers this result. The following sharp thin-shell estimate is proved in \cite{ChenKlartag}
\[
  \Var(|X|^2)\leq8n
\]
for isotropic log-concave random vectors. A minor anisotropic
adaptation of their proof gives
\begin{equation}\label{eq_401}
  \Var(|X|^2)
  \leq
  8\Tr(\Sigma^2)
\end{equation}
for every centered log-concave random vector $X$ with covariance
matrix $\Sigma$.

Consequently, if $X$ is isotropic and $A$ is symmetric positive
semidefinite, then
\[
  \Var\bigl(\langle AX,X\rangle\bigr)
  \leq
  8\Tr(A^2).
\]
For an arbitrary symmetric matrix, splitting $A$ into its positive and
negative parts gives
\[
  \Var\bigl(\langle AX,X\rangle\bigr)
  \leq
  16\Tr(A^2).
\]
This is sufficient to obtain
\begin{equation}\label{eq_424}
  \kappa_n^2\leq16,
\end{equation}
see Corollary~\ref{cor_915} below. With a little more work one may
obtain the slightly better estimate $\kappa_n^2\leq8$, but this is not
relevant for our purposes.

The parameter $\kappa_n$ first appeared in Eldan's original paper on
the stochastic-localization process~\cite{Eldan}. Let us briefly
recall its role. We refer to the lecture notes of Klartag and
Lehec~\cite{KlartagLehecSurvey} for background on the process, its
history and applications.

Let $X\sim\mu$ be isotropic and log-concave, let $(B_t)_{t\geq0}$ be an
independent Brownian motion, and let
\[
  \mu_t=\mathcal L\bigl(X\,\big|\,tX+B_t\bigr).
\]
The stochastic process
$(\mu_t)_{t\geq0}$ is called the stochastic-localization process
initiated at $\mu$. Write $A_t=\Cov(\mu_t)$, so that $A_0=\Id$. Put
\begin{equation}\label{eq_445}
  T_0=\frac{c}{\kappa_n^2\log n}
\end{equation}
for an appropriate constant $c>0$. With probability at least
$1-Ce^{-c/T_0}$, one has
\begin{equation}\label{eq_450}
  \frac12\Id\preceq A_t\preceq2\Id,
  \qquad 0\leq t\leq T_0;
\end{equation}
see~\cite[Proposition~2.3]{BizeulKlartag}. The upper bound in
\eqref{eq_450} is essentially due to Eldan~\cite{Eldan},
while the lower bound was established by Bizeul and
Klartag~\cite{BizeulKlartag}, following previous work of Lee and
Vempala~\cite{LeeVempala}.

In view of~\eqref{eq_424}, we may take $T_0=c/\log n$, and, for
a suitable choice of $c>0$, the failure probability in
\eqref{eq_450} is at most $n^{-10}$. The order $1/\log n$ is
sharp. It is explained in~\cite{KlartagLehecSurvey} that products of
centered exponentials saturate the upper bound, while an analogous
argument shows that the uniform distribution on the cube saturates the
lower bound.

In~\cite{KlartagKLS} Klartag establishes an improved log-concave Lichnerowicz inequality and uses it to prove that 
\begin{equation}\label{eq_469}
    \Psi_n^2 \lesssim T_0^{-1/2},
\end{equation}
where $\Psi_n$ is the KLS constant~\cite{KannanLovaszSimonovits}, that is, the best constant such that for every isotropic log-concave vector $X$, and every Lipschitz function $f$, we have the Poincaré inequality
$$\Var f(X) \leq \Psi_n^2 \, \E |\nabla f(X)|^2.$$
Therefore, using~\eqref{eq_445},~\eqref{eq_469} and the
dimension-free estimate~\eqref{eq_424}, one obtains
\begin{equation}\label{eq_476}
  \Psi_n\leq C(\log n)^{1/4},
\end{equation}
as recorded in~\cite{Letwin}.

The same covariance estimate underlies the gauge comparisons in
\cite{EldanLehec,BizeulKlartag}, which were established with a factor of
order $T_0^{-1/2}$. Therefore, using~\eqref{eq_445} and
\eqref{eq_424} yields Theorem~\ref{thm_206}.

\medskip

The paper is organized as follows. Section~\ref{sec_513} is purely
expository. It provides a sketch of the
argument of Chen and Klartag and explains the minor anisotropic adaptation of their proof yielding the sharp generalized variance bound, and its application to $\kappa_n$, as recorded in \cite{Letwin}. Section~\ref{sec_940} proves
Theorem~\ref{thm_230} and uses it to deduce Theorem~\ref{thm_186}. It also contains a discussion on the optimality of the results, and records applications of the $MM^*$
estimate to Banach-Mazur distances and the flatness constant.

\paragraph{Notation.}
Throughout the paper, $C,c>0$ denote universal constants whose values
may change from one occurrence to the next. We write $A\lesssim B$ if
$A\leq CB$, and $A\simeq B$ if both $A\lesssim B$ and $B\lesssim A$.
The notation $|x|$ denotes the Euclidean norm of a vector, whereas
$|K|$ denotes the Lebesgue volume of a measurable set. For a matrix
$A$, we write
\[
  \norm{A}_{\HS}=\sqrt{\Tr(A^\top A)},
\]
and $A\preceq B$ means that $B-A$ is positive semidefinite. Finally,
for $x,y\in\R^n$, we set $x\otimes y=(x_i y_j)_{i,j=1}^n$.

\paragraph{Acknowledgements} We would like to thank Boaz Klartag, Dan Mikulincer and Beatrice-Helen Vritsiou for many insightful and stimulating conversations about $MM^*$ estimate and for their feedback on an earlier version of this manuscript. We learnt about Conjecture \ref{conj_358} from Emanuel Milman.

\section{An anisotropic version of the Chen-Klartag argument}
\label{sec_513}

We stress that this section is purely expository and contains no
original content. It is based entirely on the preprint of Chen and Klartag
\cite{ChenKlartag}, and describes the minor anisotropic adaptation of
their argument, first recorded in~\cite{Letwin}, which yields a
generalized thin-shell estimate.

The thin-shell problem originates in the work of Anttila, Ball and
Perissinaki~\cite{AnttilaBallPerissinaki} and was subsequently
formulated as the variance conjecture by Bobkov and
Koldobsky~\cite{BobkovKoldobsky}. It asks whether every isotropic
log-concave random vector $X$ in $\R^n$ satisfies
\[
  \Var(|X|^2)\leq Cn.
\]
After a long sequence of developments, this conjecture was recently
resolved by Klartag and Lehec~\cite{KlartagLehec}, following a
breakthrough of Guan~\cite{Guan}. Most of the preceding approaches
used Eldan's stochastic-localization process; see, among others,
\cite{Eldan,LeeVempala,KlartagLehecPolylog,KlartagKLS,Guan,KlartagLehec}.
Chen and Klartag~\cite{ChenKlartag} gave a different proof, based on
log-concave moment measures, and obtained the optimal constant $8$.
The purpose of this section is to explain the minor anisotropic
modification of their argument which gives
\begin{equation}\label{eq_541}
  \Var(|X|^2)\leq8\Tr(\Sigma^2)
\end{equation}
for a centered log-concave vector with covariance matrix $\Sigma$, as
recorded in~\cite{Letwin}. The question of studying thin-shell
estimates beyond the isotropic setting was put forward
in~\cite{AlonsoBastero}.

A common starting point in recent approaches to thin-shell problems is
the $H^{-1}$ inequality of Klartag~\cite{KlartagHminus}, which was extended to general log-concave measures in~\cite{BartheKlartag}. For a centered function
$f\in L^2(\mu)$, set
\begin{equation}\label{eq_554}
  \norm{f}_{H^{-1}(\mu)}
  =
  \sup\left\{
    \int fg\dd\mu:
    g\in C_c^\infty(\R^n),\
    \int|\nabla g|^2\dd\mu\leq1
  \right\}.
\end{equation}
Here and below, $ C_c^\infty(\R^n)$ is the class of smooth compactly supported functions. The $H^{-1}$ inequality asserts that if $f$ is locally Lipschitz and 
$f,\partial_1f,\ldots,\partial_nf\in L^2(\mu)$ are centered,
then
\begin{equation}\label{eq_570}
  \Var_\mu(f)
  \leq
  \sum_{i=1}^n\norm{\partial_i f}_{H^{-1}(\mu)}^2.
\end{equation}
Applying this inequality to
$f(x)=|x|^2-\int|x|^2\dd\mu(x)$ gives
\begin{equation}\label{eq_577}
  \Var(|X|^2)
  \leq
  4\sum_{i=1}^n\norm{x_i}_{H^{-1}(\mu)}^2.
\end{equation}

The right-hand side admits a simple interpretation in terms of Stein
kernels. Recall that a square-integrable matrix field
$\tau:\R^n\to\R^{n\times n}$ is a Stein kernel for $\mu$ if
\begin{equation}\label{eq_586}
  \int x\cdot F(x)\dd\mu(x)
  =
  \int\ip{\tau(x)}{DF(x)}_{\HS}\dd\mu(x)
\end{equation}
for every $F\in C_c^\infty(\R^n;\R^n)$, where
$DF$ denotes the Jacobian matrix of $F$.

\begin{proposition}\label{prop_594}
Let $\tau$ be a Stein kernel for $\mu$. Then
\begin{equation}\label{eq_596}
  \Var(|X|^2)
  \leq
  4\int\norm{\tau(x)}_{\HS}^2\dd\mu(x).
\end{equation}
\end{proposition}

\begin{proof}
For $g\in C_c^\infty(\R^n)$, applying the Stein identity to
$F=ge_i$ gives
\[
  \int x_i g\dd\mu
  =
  \int\sum_{j=1}^n\tau_{ij}\partial_jg\dd\mu.
\]
By Cauchy-Schwarz,
\[
  \norm{x_i}_{H^{-1}(\mu)}^2
  \leq
  \int\sum_{j=1}^n\tau_{ij}^2\dd\mu.
\]
Summing over $i$ and using~\eqref{eq_577}
proves~\eqref{eq_596}.
\end{proof}

In fact, Courtade, Fathi and Pananjady
\cite{CourtadeFathiPananjady} construct a distinguished Stein kernel
$\tauL$, related to the Langevin generator of $\mu$, satisfying
\begin{equation}\label{eq_624}
  \int\norm{\tauL}_{\HS}^2\dd\mu
  =
  \sum_{i=1}^n\norm{x_i}_{H^{-1}(\mu)}^2.
\end{equation}
Thus the thin-shell problem is reduced to constructing a Stein kernel
with sufficiently small Hilbert-Schmidt energy. The proof of Chen and
Klartag uses a different kernel, built from moment measures, which we
now describe.

\subsection{Moment measures and the Chen-Klartag argument}

Let $\mu$ be a centered log-concave probability measure on $\R^n$,
with covariance matrix
\[
  \Sigma=\int x\otimes x\dd\mu(x).
\]
The moment-measure theorem of Cordero-Erausquin and
Klartag~\cite{CorderoKlartag} provides an essentially continuous
convex function $\psi:\R^n\to\R\cup\{+\infty\}$, unique up to
translation of its argument, such that
\begin{equation}\label{eq_645}
  \int_{\R^n}e^{-\psi(y)}\dd y=1,
  \qquad
  (\nabla\psi)_\#\bigl(e^{-\psi(y)}\dd y\bigr)=\mu.
\end{equation}

Following Chen and Klartag~\cite{ChenKlartag}, we assume temporarily
that $\mu$ has density $e^{-V}$ on a bounded open convex set $K$, with
smooth boundary, and that $V$ is smooth and convex in a neighborhood
of $K$. Under these assumptions, Klartag~\cite{KlartagMoment} proved
that $\psi$ is smooth and strictly convex, and that
$\nabla\psi:\R^n\to K$ is a diffeomorphism. Set
\begin{equation}\label{eq_657}
  \dd\nu(y)=e^{-\psi(y)}\dd y,
  \qquad
  H(y)=D^2\psi(y).
\end{equation}

The Hessian $H$, transported to the target measure, is a symmetric
positive-definite Stein kernel, as observed by
Fathi~\cite[Theorem~2.3]{Fathi}. We denote it by
\begin{equation}\label{eq_666}
  \tauM(x)=H\bigl((\nabla\psi)^{-1}(x)\bigr).
\end{equation}
Indeed, formally, if $F\in C_c^\infty(\R^n;\R^n)$, integration by
parts and the push-forward relation give
\begin{align*}
  \int_Kx\cdot F(x)\dd\mu(x)
  &=\int_{\R^n}\nabla\psi(y)\cdot F(\nabla\psi(y))\dd\nu(y)\\
  &=\int_{\R^n}\Tr\bigl(H(y)DF(\nabla\psi(y))\bigr)\dd\nu(y)\\
  &=\int_K\ip{\tauM(x)}{DF(x)}_{\HS}\dd\mu(x).
\end{align*}
Moreover,
\begin{equation}\label{eq_678}
  \int_{\R^n}H\dd\nu
  =\int_Kx\otimes x\dd\mu(x)
  =\Sigma.
\end{equation}

The change-of-variables formula in~\eqref{eq_645} yields the
Monge-Amp\`ere equation
\begin{equation}\label{eq_686}
  e^{-V(\nabla\psi)}\det H=e^{-\psi}.
\end{equation}
As explained in \cite{KlartagMoment}, it is useful to regard $H$ as a Riemannian
metric on $\R^n$, endowed with the probability measure
$\nu=e^{-\psi}\dd y$. The associated symmetric generator is
\begin{equation}\label{eq_692}
  \cL f
  =e^\psi\operatorname{div}\bigl(e^{-\psi}H^{-1}\nabla f\bigr),
\end{equation}
and its carr\'e du champ is
\begin{equation}\label{eq_697}
  \Gamma(f,g)=\ip{H^{-1}\nabla f}{\nabla g}.
\end{equation}
Thus
\[
  \int g\,\cL f\dd\nu=-\int\Gamma(f,g)\dd\nu.
\]
The Brascamp-Lieb inequality~\cite{BrascampLieb} for the measure
$\nu$ reads
\begin{equation}\label{eq_706}
  \Var_\nu(f)\leq\int\Gamma(f)\dd\nu,
\end{equation}
where $\Gamma(f)=\Gamma(f,f)$. When $\cL$ is applied to a matrix
field, it acts entrywise.

We now give a brief overview of the part of the Chen-Klartag argument
used below. Write
\[
  \psi_i=\partial_i\psi,
  \qquad
  \psi_{ij}=\partial_i\partial_j\psi,
  \qquad
  \psi_{ijk}=\partial_i\partial_j\partial_k\psi.
\]
For $i=1,\ldots,n$, put
\[
  C_i=\partial_iH,
  \qquad (C_i)_{jk}=\psi_{ijk}.
\]
Differentiating the Monge-Amp\`ere equation produces the
nonnegative matrix fields
\begin{equation}\label{eq_728}
  A=H\bigl(D^2V\circ\nabla\psi\bigr)H\succeq0,
  \qquad
  Q_{ij}=\Tr(H^{-1}C_iH^{-1}C_j),
  \quad Q\succeq0,
\end{equation}
and the identity
\begin{equation}\label{eq_735}
  \cL H+H=A+Q.
\end{equation}
This is Lemma~3.1 of~\cite{ChenKlartag}. It refines an earlier argument of Klartag \cite{KlartagMoment} who established the positivity of the left-hand side of \eqref{eq_735}.

We are interested in bounding the Hilbert-Schmidt norm of the kernel
$\tauM$, namely
\begin{equation}\label{eq_742}
  N_2=\int_{\R^n}\Tr(H^2)\dd\nu
  =\int\norm{\tauM(x)}_{\HS}^2\dd\mu(x).
\end{equation}
At a fixed point, choose coordinates in which
\[
  H=\operatorname{diag}(\lambda_1,\ldots,\lambda_n),
  \qquad \lambda_i>0,
\]
and define
\begin{align*}
  d_2
  &=\sum_{a,i,j=1}^n\frac{\psi_{aij}^2}{\lambda_a},\\
  a_2&=\Tr(HA),\\
  q_2
  &=\Tr(HQ)
  =\sum_{a,i,j=1}^n
  \frac{\lambda_a}{\lambda_i\lambda_j}\psi_{aij}^2.
\end{align*}
Put
\begin{equation}\label{eq_762}
  D_2=\int d_2\dd\nu,
  \qquad
  A_2=\int a_2\dd\nu,
  \qquad
  Q_2=\int q_2\dd\nu.
\end{equation}
The two identities used below are Lemmas~3.3 and~3.5
of~\cite{ChenKlartag}.

\begin{lemma}[Chen-Klartag]\label{lem_772}
Under the preceding regularity assumptions,
\begin{equation}\label{eq_774}
  N_2=D_2+A_2+Q_2,
\end{equation}
and
\begin{equation}\label{eq_778}
  q_2-d_2
  =\frac16\sum_{a,i,j=1}^n
  \frac{\psi_{aij}^2}{\lambda_a\lambda_i\lambda_j}
  \Bigl[(\lambda_a-\lambda_i)^2
       +(\lambda_i-\lambda_j)^2
       +(\lambda_j-\lambda_a)^2\Bigr]\geq0.
\end{equation}
In particular, $Q_2\geq D_2$.
\end{lemma}

The identity~\eqref{eq_774} follows by applying the
diffusion product rule to $\Tr(H^2)$ and integrating~\eqref{eq_735}.
The inequality~\eqref{eq_778} relies on the symmetry of the third
derivative tensor $D^3\psi$. We may now state the Chen-Klartag
bootstrap~\cite[Theorem~1.5]{ChenKlartag}.

\begin{theorem}[Chen-Klartag]\label{thm_795}
If $\mu$ is isotropic, then
\begin{equation}\label{eq_797}
  \int\norm{\tauM(x)}_{\HS}^2\dd\mu(x) = \int_{\R^n}\Tr(H^2)\dd\nu\leq2n.
\end{equation}
\end{theorem}

\begin{proof}
If $\mu$ is isotropic, then~\eqref{eq_678} gives
$\int H\dd\nu=\Id$. Applying~\eqref{eq_706} to every entry of
$H$ yields
\begin{equation}\label{eq_806}
  N_2-n
  =\int\norm{H-\Id}_{\HS}^2\dd\nu
  \leq D_2.
\end{equation}
On the other hand, Lemma~\ref{lem_772} gives
\[
  N_2=D_2+A_2+Q_2\geq2D_2.
\]
Combining the two inequalities gives $N_2\leq2n$.
\end{proof}

\subsection{The anisotropic adaptation}

The preceding proof does not use isotropy until the Brascamp-Lieb
step~\eqref{eq_806}. Replacing that single line
gives the following anisotropic version.

\begin{proposition}\label{prop_824}
If $\mu$ has covariance matrix $\Sigma$,
\begin{equation}\label{eq_826}
  \int\norm{\tauM(x)}_{\HS}^2\dd\mu(x)
  =\int_{\R^n}\Tr(H^2)\dd\nu
  \leq2\Tr(\Sigma^2).
\end{equation}
\end{proposition}

\begin{proof}
By~\eqref{eq_678}, $\int H\dd\nu=\Sigma$. Applying the
Brascamp-Lieb inequality~\eqref{eq_706} to every entry of $H$
and summing gives
\begin{align}
  N_2-\Tr(\Sigma^2)
  &=\int\norm{H-\Sigma}_{\HS}^2\dd\nu\notag\\
  &=\sum_{i,j=1}^n\Var_\nu(H_{ij})
  \leq\sum_{i,j=1}^n\int\Gamma(H_{ij})\dd\nu
  =D_2.
  \label{eq_843}
\end{align}
Since $A_2\geq0$ and $Q_2\geq D_2$, we still have
\[
  N_2=D_2+A_2+Q_2\geq2D_2.
\]
Using~\eqref{eq_843}, we conclude that
\[
  N_2\geq2\bigl(N_2-\Tr(\Sigma^2)\bigr),
\]
which is equivalent to~\eqref{eq_826}.
\end{proof}

Therefore, we arrive at the following sharp anisotropic thin-shell
estimate, recorded in~\cite{Letwin}.

\begin{theorem}\label{thm_859}
Let $X$ be a centered log-concave random vector in $\R^n$ with
covariance matrix $\Sigma$. Then
\begin{equation}\label{eq_862}
  \Var(|X|^2)\leq8\Tr(\Sigma^2).
\end{equation}
The constant $8$ is optimal.
\end{theorem}

\begin{proof}
Under the regularity assumptions, Proposition~\ref{prop_594}
applied to the Stein kernel $\tauM$ and
Proposition~\ref{prop_824} give
\[
  \Var(|X|^2)
  \leq4\int\norm{\tauM}_{\HS}^2\dd\mu
  \leq8\Tr(\Sigma^2).
\]
As explained in~\cite[Section~3]{ChenKlartag}, the regularity
assumptions may be removed by approximation. Sharpness follows by
considering independent centered exponential variables with the
appropriate variances in a basis diagonalizing $\Sigma$.
\end{proof}

We next pass from~\eqref{eq_862} to quadratic forms.

\begin{corollary}\label{cor_885}
Let $X$ be an isotropic log-concave random vector in $\R^n$. If
$A\succeq0$ is symmetric, then
\begin{equation}\label{eq_888}
  \Var\bigl(\langle AX,X\rangle\bigr)
  \leq8\Tr(A^2).
\end{equation}
For an arbitrary symmetric matrix $A$,
\begin{equation}\label{eq_893}
  \Var\bigl(\langle AX,X\rangle\bigr)
  \leq16\Tr(A^2).
\end{equation}
\end{corollary}

\begin{proof}
For the first statement, apply Theorem~\ref{thm_859}
to $Y=A^{1/2}X$. For a general symmetric matrix, split it into positive and negative parts
$A=A_+-A_-$, where $A_+,A_-\succeq0$ and $A_+A_-=0$ and use 
\eqref{eq_888}.
\begin{align*}
  \Var\bigl(\langle AX,X\rangle\bigr)
  &\leq16\Tr(A_+^2)+16\Tr(A_-^2)\\
  &=16\Tr(A^2).
\end{align*}
\end{proof}

We now return to the third-moment parameter $\kappa_n$ defined
in~\eqref{eq_381}.

\begin{corollary}\label{cor_915}
One has
\begin{equation}\label{eq_917}
  \kappa_n^2\leq16.
\end{equation}
\end{corollary}

\begin{proof}
Let $X$ be isotropic and log-concave, let
$\theta\in\mathbb S^{n-1}$, and set
\[
  B=\E\bigl[\langle X,\theta\rangle X\otimes X\bigr].
\]
Then $B$ is symmetric and
\begin{align*}
  \norm{B}_{\HS}^2
  &=\E\left(\langle X,\theta\rangle.
               \langle BX,X\rangle\right)\\
  &\leq \Var\bigl(\langle BX,X\rangle\bigr)^{1/2}
  \leq4\norm{B}_{\HS},
\end{align*}
where we used isotropy and~\eqref{eq_893}. Hence
$\norm{B}_{\HS}\leq4$, and the result follows.
\end{proof}

\section{Mean-width bound and applications}\label{sec_940}

In this section we prove the main result of the paper. We first prove the
geometric inequality of Theorem~\ref{thm_230} using heat flow and combine it with the
gauge-order comparison of Theorem~\ref{thm_206} to deduce
Theorem~\ref{thm_186}. We then establish the isotropic and affine optimality statements from the introduction, before recording the
consequences for Banach-Mazur distances and the flatness constant.

Before anything, we recall the very standard fact that spherical integration may be replaced by Gaussian integrations in the definition of the mean-norm and mean-width. Indeed if $G$ is a standard Gaussian vector in $\R^n$, then $G = |G|\Theta$, where $\Theta$ is uniformly distributed on the sphere, $|G|$ is the Euclidean norm of $G$ and both quantities are independent. Therefore, by homogeneity, for any $K$,
\begin{equation}\label{eq_949}
  \E\norm{G}_K=\E|G|\,M(K) \simeq \sqrt{n}M(K),
  \qquad
  \E h_K(G)=\E|G|\,M^*(K) \simeq \sqrt{n}M^*(K),
\end{equation}
where we used that $\E|G|\simeq\sqrt n.$

\subsection{A heat-flow argument}

The main new ingredient for the mean-width estimate is
Theorem~\ref{thm_230}, whose proof we now give.

\begin{proof}[Proof of Theorem~\ref{thm_230}]
By a standard approximation argument, we may assume that $K$ has a smooth boundary. Let $X_K$ be uniformly distributed in the centered
convex body $K$, and write
\[
  W_K=\E h_K(X_K),
  \qquad
  W_G=\E h_K(G).
\]
Let $(B_s)_{s\geq0}$ be a standard Brownian motion, independent of
$X_K$, and set
\[
  F(s)=\E h_K(X_K+B_s)
  =\frac{1}{|K|}\int_K P_s h_K(x)\dd x,
\]
where $P_s$ is the heat semigroup with generator $\Delta/2$. By the
divergence theorem,
\[
  F'(s)
  =\frac{1}{2|K|}
  \int_{\partial K}
  \ip{\nabla P_s h_K(x)}{n_K(x)}\dd S(x),
\]
where $n_K$ denotes the outer unit normal and $\dd S$ the surface measure.
The support function is differentiable almost everywhere and
$\nabla h_K(y)\in K$ at every differentiability point. Hence, by convexity of $K$, the average
\[
  \nabla P_s h_K=P_s(\nabla h_K)
\]
also takes values in $K$. For $x\in\partial K$,
\[
  \ip{\nabla P_s h_K(x)}{n_K(x)}
  \leq h_K(n_K(x))
  =\ip{x}{n_K(x)}.
\]
A second application of the divergence theorem gives
\[
  F'(s)
  \leq\frac{1}{2|K|}
  \int_{\partial K}\ip{x}{n_K(x)}\dd S(x)
  =\frac n2.
\]
Consequently,
\[
  F(s)\leq W_K+\frac{ns}{2}.
\]
Choosing $s_*=2W_K/n$, we find that $F(s_*)\leq2W_K$. On the other hand, since $X_K$ is centered, Jensen's inequality yields
\[
  F(s_*)
  =\E h_K(X_K+\sqrt{s_*}G)
  \geq\E h_K(\sqrt{s_*}G)
  =\sqrt{s_*}\,W_G.
\]
Therefore, combining both inequalities yields
\[
  W_G^2\leq\frac{4W_K^2}{s_*}=2nW_K,
\]
and the proof is complete.
\end{proof}

\subsection{Proof of Theorem~\ref{thm_186}}

We now deduce Theorem~\ref{thm_186}. Let $K$ be isotropic. Since
$h_K=\norm{\cdot}_{K^\circ}$ is a gauge,
Theorem~\ref{thm_206} gives
\[
  \E h_K(X_K)
  \leq C\sqrt{\log n}\,\E h_K(G).
\]
Combining this estimate with Theorem~\ref{thm_230}, we obtain
\[
  \bigl(\E h_K(G)\bigr)^2
  \leq Cn\sqrt{\log n}\,\E h_K(G),
\]
and hence
\[
  \E h_K(G)\leq Cn\sqrt{\log n}.
\]
Using~\eqref{eq_949}, we conclude that
\[
  M^*(K)\leq C\sqrt{n\log n},
\]
as claimed.

\subsection{Optimality of the estimates}\label{sec_1045}

We now prove the optimality statements made in the introduction. We first
consider the isotropic position and then pass to arbitrary affine positions.

\paragraph{Sharpness in isotropic position.}
Consider the convex bodies
\[
  Q_n=\sqrt3\,B_\infty^n,
  \qquad
  P_n=\sqrt{\frac{(n+1)(n+2)}2}\,B_1^n.
\]
The bodies $Q_n$ and $P_n$ are respectively the isotropic cube and the
isotropic cross-polytope. Since
$\E\norm{G}_\infty\simeq\sqrt{\log n}$, formula
\eqref{eq_949} gives
\begin{equation}\label{eq_1061}
  M(Q_n)
  \simeq\sqrt{\frac{\log n}{n}},
  \qquad
  M^*(P_n)
  \simeq\sqrt{n\log n}.
\end{equation}
Thus the mean-norm and mean-width estimates are separately sharp in
isotropic position, even among origin-symmetric bodies.

The product estimate is also sharp in the symmetric isotropic class. For
simplicity, let the ambient dimension be $2m$ and consider
\[
  K=Q_m\times P_m.
\]
The body $K$ is origin-symmetric and isotropic. Let
$G=(G_1,G_2)$, where $G_1,G_2$ are independent standard Gaussian vectors
in $\R^m$. We have
\[
  \norm{(x,y)}_K
  =\max\{\norm{x}_{Q_m},\norm{y}_{P_m}\},
  \qquad
  h_K(x,y)=h_{Q_m}(x)+h_{P_m}(y).
\]
Moreover,
\[
  \E\norm{G_1}_{Q_m}\simeq\sqrt{\log m},
  \qquad
  \E\norm{G_2}_{P_m}\simeq1,
\]
and
\[
  \E h_{Q_m}(G_1)\simeq m,
  \qquad
  \E h_{P_m}(G_2)\simeq m\sqrt{\log m}.
\]
Therefore,
\[
  M(K)\simeq\sqrt{\frac{\log m}{m}},
  \qquad
  M^*(K)\simeq\sqrt{m\log m},
\]
and hence
\begin{equation}\label{eq_1104}
  \ell(K)\simeq\log m.
\end{equation}
This proves the sharpness of the isotropic $MM^*$ estimate, even in the
origin-symmetric class.

This example is nevertheless compatible with
Conjecture~\ref{conj_358}. Indeed, for the linear image
\[
  \widetilde K
  =Q_m\times\frac1{\sqrt{\log m}}P_m,
\]
the same computation gives
\[
  M(\widetilde K)\simeq\sqrt{\frac{\log m}{m}},
  \qquad
  M^*(\widetilde K)\simeq\sqrt m,
  \qquad
  \ell(\widetilde K)\simeq\sqrt{\log m}.
\]

\paragraph{Minimal mean-width position.}
We shall use a characterization of the minimal mean-width position due to Giannopoulos and Milman~\cite{GiannopoulosMilman}.
Let $\sigma$ denote the rotationally invariant probability measure on
$\mathbb S^{n-1}$. We say that a convex body $K$ is in minimal mean-width position if
\[
  M^*(K)\leq M^*(TK)
  \qquad\text{for every }T\in SL_n.
\]
Giannopoulos and Milman proved that this is equivalent to the measure $h_K\dd\sigma$ having a scalar covariance matrix,
\begin{equation}\label{eq_1135}
  \int_{\mathbb S^{n-1}}h_K(u)\,u\otimes u\dd\sigma(u)
  =\frac{M^*(K)}{n}\,\Id.
\end{equation}

If $K$ is $1$-symmetric in some basis, that is, invariant under coordinate permutations and sign changes, so is the measure $h_K\dd\sigma$, therefore its covariance matrix is scalar and $K$ is in minimal mean-width position.

The same conclusion applies to the mean-norm $M$ : If $K$ is $1$-symmetric then it is in minimal mean-norm position. Indeed, by what precedes, its polar $K^\circ$, which inherits the same symmetries, has minimal mean width. Now for any $T\in SL_n$,
\[
  M(TK)=M^*((TK)^\circ)
  =M^*((T^{-1})^TK^\circ)
  \geq M^*(K^\circ)=M(K).
\]
It follows that the standard positions of the cube and the cross-polytope
simultaneously minimize $M$ and $M^*$ over volume-preserving linear images.

When working in the symmetric class, it is enough to consider linear images of $K$, rather than affine ones. Indeed, while the mean-width is invariant by translation, the mean-norm can only increase if an origin-symmetric body is translated. For completeness, we give a short proof of this elementary fact.

For a convex body $L$ containing the origin, write
\[
  \rho_L(u)=\max\{t\geq0:\ tu\in L\},
\]
so that $\norm{u}_L=1/\rho_L(u)$. Let $K=-K$ and assume that
$0\in\operatorname{int}(K+b)$. Put
$r_+=\rho_{K+b}(u)$ and $r_-=\rho_{K+b}(-u)$. Then
$r_+u-b\in K$, while symmetry gives $r_-u+b\in K$. By convexity,
\[
  \frac{r_++r_-}{2}u\in K,
\]
and therefore $r_++r_-\leq2\rho_K(u)$. Hence
\[
  \frac12\bigl(\norm{u}_{K+b}+\norm{-u}_{K+b}\bigr)
  =\frac12\left(\frac1{r_+}+\frac1{r_-}\right)
  \geq\frac2{r_++r_-}
  \geq\norm{u}_K.
\]
Integrating over the sphere yields
\begin{equation}\label{eq_1172}
  M(K+b)\geq M(K).
\end{equation}

\paragraph{Affine optimality of $M$ and $M^*$.}
We now prove Corollary~\ref{cor_319}. Let $K$ be a convex body
with $|K|=|B_2^n|$ and let $\widetilde K = TK + a$, $T\in SL_n$ be an affine image of $K$ such that $\widetilde K$ is centered and
$$\Cov(\widetilde K) = \frac{1}{|\widetilde K|}\int_{\widetilde K}x\otimes x \,\dd x = \lambda_K^{-2} I_n$$
for some $\lambda_K>0$. The solution of the slicing
problem~\cite{KlartagLehecSlicing, BizeulSlicing, BrazitikosSlicing} says exactly that 
$$\lambda_K \simeq \sqrt{n}.$$
Applying Theorems~\ref{thm_252} and~\ref{thm_186} to the isotropic convex body $\lambda_K \widetilde K$ , we obtain
\[
  M(\widetilde K)\lesssim\sqrt{\log n},
  \qquad
  M^*(\widetilde K)\lesssim\sqrt{\log n}.
\]
This proves the upper bounds in~\eqref{eq_321}.

For the reverse direction, set $\widetilde P_n$ and $\widetilde Q_n$ to be scaled versions of $Q_n$ and $P_n$ so that their volumes
are equal to $|B_2^n|$. The scaling factors are of order $n^{-1/2}$, so~\eqref{eq_1061} becomes
\[
  M(\widetilde Q_n)\simeq\sqrt{\log n},
  \qquad
  M^*(\widetilde P_n)\simeq\sqrt{\log n}.
\]
By the preceding paragraph, $\widetilde Q_n$ and $\widetilde P_n$ are in minimal mean-norm and mean-width position. This proves the corollary, including
its restriction to origin-symmetric convex bodies.

\begin{remark}
    In the affine optimality argument above, the resolution of the slicing problem is only needed to show that the $M^*$ bound in isotropic position yield an affine optimal bound. For the $M$ bound, we only need the "easy" direction of the slicing problem, namely that any convex body in isotropic position has a volume radius at most $\sqrt{n}$, which can be proved in an elementary fashion. 

    Therefore, as discussed in the introduction, one may dualize the $M$ bound to prove get an affine optimal $M^*$ bound when $K$ is in "dual isotropic position", that is when $K^\circ$ is isotropic. This latter argument does not require a solution to the slicing problem.
\end{remark}

\paragraph{The affine $MM^*$ problem.}
We finally prove Corollary~\ref{cor_340}. The upper bound in
\eqref{eq_342} follows by placing an arbitrary convex body in
isotropic position and applying Corollary~\ref{cor_265}.

The reverse direction was settled by Banaszczyk, Litvak, Pajor and
Szarek  \cite{BanaszczykLitvakPajorSzarek}, who proved that for any simplex $\Delta$ containing the origin in its interior,
\[
  \ell(\Delta)\gtrsim\log n.
\]
Together with the upper bound, this proves
\eqref{eq_342}.

Finally, the conjectural $\sqrt{\log n}$ order in the origin-symmetric class is itself
optimal. Indeed, as previously discussed, the standard position of the cube $B_\infty^n$ (or equivalently, the cross polytope) simultaneously
minimizes $M$ and $M^*$ over $SL_n$. Since $\ell$ is scale invariant, this
implies
\[
  \inf_{T\in GL_n}\ell(TB_\infty^n)
  =\ell(B_\infty^n)
  \simeq\sqrt{\log n}.
\]
On the other hand,~\eqref{eq_1104} shows that the
isotropic position may have an $MM^*$ product of order $\log n$, even in the
origin-symmetric class. Therefore, a proof of Conjecture \ref{conj_358} should not use the isotropic position.

\subsection{Applications of the MM* estimate}

We conclude with two consequences of the bounds on $M$ and $M^*$.
Recall that the Banach-Mazur distance between two convex bodies
$K,L\subset\R^n$ is defined by
\[
  d_{BM}(K,L)
  =\inf\left\{\lambda\geq1:\ \exists T\in GL_n,\ x,y\in\R^n,
  \ TK+x\subset L+y\subset\lambda(TK+x)\right\}.
\]
The argument of~\cite{BizeulKlartag} gives, after putting
$K,L\subset\R^n$ in isotropic position,
\[
  d_{BM}(K,L)
  \leq C\Bigl(
  n\max\{M(K),M(L)\}
  +\max\{M^*(K),M^*(L)\}
  \Bigr)^2
\]
for all $n\geq2$. Together with~\eqref{eq_188} and
\eqref{eq_254}, this yields
\begin{equation}\label{eq_1250}
  d_{BM}(K,L)\leq Cn\log n.
\end{equation}
The relevance of $MM^*$ bounds to Banach-Mazur distances was pioneered by Rudelson \cite{Rudelson}.

The same estimates also improve the flatness bound. Recall that the
flatness constant $\operatorname{Flt}(n)$ is the smallest $F$ such
that every convex body $K\subset\R^n$ satisfying
$K\cap\Z^n=\varnothing$ has lattice width at most $F$, that is,
\[
  \min_{z\in\Z^n\setminus\{0\}}
  \left(
  \sup_{x\in K}\langle x,z\rangle
  -\inf_{x\in K}\langle x,z\rangle
  \right)
  \leq F.
\]
It follows from the $MM^*$ estimate and the argument of Banaszczyk,
Litvak, Pajor and Szarek~\cite{BanaszczykLitvakPajorSzarek} that
\begin{equation}\label{eq_1268}
  \operatorname{Flt}(n)\leq Cn\log n.
\end{equation}
This improves the bound $O(n\log^2 n)$ obtained by Reis and
Rothvoss~\cite{ReisRothvoss}.




\begin{thebibliography}{99}

\bibitem{AlonsoBastero}
D.~Alonso-Guti\'errez and J.~Bastero,
\emph{The variance conjecture on some polytopes},
in Asymptotic Geometric Analysis, Fields Inst. Commun. 68,
Springer, New York, 2013, 1-20.

\bibitem{AnttilaBallPerissinaki}
M.~Anttila, K.~Ball, and I.~Perissinaki,
\emph{The central limit problem for convex bodies},
Trans. Amer. Math. Soc. \textbf{355} (2003), no.~12, 4723-4735.

\bibitem{ArtsteinAvidanGiannopoulosMilman}
S.~Artstein-Avidan, A.~Giannopoulos, and V.~D.~Milman,
\emph{Asymptotic Geometric Analysis, Part I},
Mathematical Surveys and Monographs, Vol.~202,
American Mathematical Society, Providence, RI, 2015.

% \bibitem{AubrunSzarek}
% G.~Aubrun and S.~J.~Szarek,
% \emph{Alice and Bob Meet Banach: The Interface of Asymptotic Geometric Analysis and Quantum Information Theory},
% Mathematical Surveys and Monographs, Vol.~223,
% American Mathematical Society, Providence, RI, 2017.

\bibitem{BanaszczykLitvakPajorSzarek}
W.~Banaszczyk, A.~E.~Litvak, A.~Pajor, and S.~J.~Szarek,
\emph{The flatness theorem for nonsymmetric convex bodies via the local theory of Banach spaces},
Math. Oper. Res. \textbf{24} (1999), no.~3, 728-750.

\bibitem{BartheKlartag}
F.~Barthe and B.~Klartag,
\emph{Spectral gaps, symmetries and log-concave perturbations},
Bull. Hellenic Math. Soc. \textbf{64} (2020), 1-31.

\bibitem{BizeulSlicing}
P.~Bizeul,
\emph{The slicing conjecture via small ball estimates},
arXiv:2501.06854, 2025.

\bibitem{BizeulKlartag}
P.~Bizeul and B.~Klartag,
\emph{Distances between non-symmetric convex bodies: optimal bounds up to polylog},
arXiv:2510.20511v3, 2026.

\bibitem{BobkovKoldobsky}
S.~G.~Bobkov and A.~Koldobsky,
\emph{On the central limit property of convex bodies},
in Geometric Aspects of Functional Analysis (2001-02),
Lecture Notes in Math., Vol.~1807, Springer, 2003, 44-52.

\bibitem{BobkovMPosition}

S.~G.~Bobkov,
\emph{On Milman's ellipsoids and $M$-position of convex bodies},
in Concentration, Functional Inequalities and Isoperimetry,
Contemp. Math. \textbf{545}, Amer. Math. Soc., Providence, RI, 2011,
23--33.

\bibitem{BourgainMilman}
J.~Bourgain and V.~D.~Milman,
\emph{New volume ratio properties for convex symmetric bodies in
$\mathbb R^n$},
Invent. Math. \textbf{88} (1987), no.~2, 319--340.


\bibitem{BrascampLieb}
H.~J.~Brascamp and E.~H.~Lieb,
\emph{On extensions of the Brunn-Minkowski and Pr\'ekopa-Leindler theorems, including inequalities for log-concave functions},
J. Funct. Anal. \textbf{22} (1976), no.~4, 366-389.

\bibitem{BrazitikosSlicing}
S.~Brazitikos,
\newblock A deterministic proof of the slicing theorem,
\newblock \emph{arXiv preprint arXiv:2608.14342}, 2026.

\bibitem{BrazitikosGiannopoulosValettasVritsiou}
S.~Brazitikos, A.~Giannopoulos, P.~Valettas, and B.-H.~Vritsiou,
\emph{Geometry of Isotropic Convex Bodies},
Mathematical Surveys and Monographs, Vol.~196,
American Mathematical Society, Providence, RI, 2014.

\bibitem{BrazitikosPandis}
S.~Brazitikos and C.~Pandis,
\emph{Geometric bounds for the mean gauge and the mean width in
isotropic position},
arXiv:2608.07216, 2026.


\bibitem{ChenKlartag}
Y.~Chen and B.~Klartag,
\emph{Digesting the proof of the sharp thin-shell inequality},
arXiv:2607.23307, 2026.

\bibitem{CorderoKlartag}
D.~Cordero-Erausquin and B.~Klartag,
\emph{Moment measures},
J. Funct. Anal. \textbf{268} (2015), no.~12, 3834-3866.

\bibitem{CourtadeFathiPananjady}
T.~A.~Courtade, M.~Fathi, and A.~Pananjady,
\emph{Existence of Stein kernels under a spectral gap, and discrepancy bounds},
Ann. Inst. Henri Poincar\'e Probab. Stat. \textbf{55} (2019), no.~2, 777-790.

\bibitem{Eldan}
R.~Eldan,
\emph{Thin shell implies spectral gap up to polylog via a stochastic localization scheme},
Geom. Funct. Anal. \textbf{23} (2013), no.~2, 532-569.

\bibitem{EldanLehec}
R.~Eldan and J.~Lehec,
\emph{Bounding the norm of a log-concave vector via thin-shell estimates},
in Geometric Aspects of Functional Analysis: Israel Seminar (GAFA) 2011-2013,
Springer, 2014, 107-122.

\bibitem{Fathi}
M.~Fathi,
\emph{Stein kernels and moment maps},
Ann. Probab. \textbf{47} (2019), no.~4, 2172-2185.

\bibitem{GiannopoulosMilman}
A.~A.~Giannopoulos and V.~D.~Milman,
\emph{Extremal problems and isotropic positions of convex bodies},
Israel J. Math. \textbf{117} (2000), 29-60.

\bibitem{GiannopoulosMilman2}
A.~Giannopoulos and E.~Milman,
\emph{$M$-estimates for isotropic convex bodies and their
$L_q$-centroid bodies},
in \emph{Geometric Aspects of Functional Analysis},
Lecture Notes in Math. \textbf{2116},
Springer, Cham, 2014, 159--182.

\bibitem{GiannopoulosStavrakakisTsolomitisVritsiou}
A.~Giannopoulos, P.~Stavrakakis, A.~Tsolomitis, and B.-H.~Vritsiou,
\emph{Geometry of the $L_q$-centroid bodies of an isotropic
log-concave measure},
Trans. Amer. Math. Soc. \textbf{367} (2015), no.~7, 4569--4593.

\bibitem{Guan}
Q.~Guan,
\emph{A note on Bourgain's slicing problem},
arXiv:2412.09075, 2024.

\bibitem{KannanLovaszSimonovits}
R.~Kannan, L.~Lov\'asz, and M.~Simonovits,
\emph{Isoperimetric problems for convex bodies and a localization lemma},
Discrete Comput. Geom. \textbf{13} (1995), 541-559.

\bibitem{KlartagHminus}
B.~Klartag,
\emph{A Berry-Esseen type inequality for convex bodies with an unconditional basis},
Probab. Theory Related Fields \textbf{145} (2009), no.~1-2, 1-33.

\bibitem{KlartagKLS}
B.~Klartag,
\emph{Logarithmic bounds for isoperimetry and slices of convex sets},
Ars Inveniendi Analytica, Paper No.~4 (2023), 1-17.

\bibitem{KlartagMoment}
B.~Klartag,
\emph{Logarithmically-concave moment measures I},
in Geometric Aspects of Functional Analysis,
Lecture Notes in Math., Vol.~2116, Springer, Cham, 2014, 231-260.

\bibitem{KlartagLehecSlicing}
B.~Klartag and J.~Lehec,
\emph{Affirmative resolution of Bourgain's slicing problem using Guan's bound},
Geom. Funct. Anal. \textbf{35} (2025), 1147-1168.

\bibitem{KlartagLehec}
B.~Klartag and J.~Lehec,
\emph{Thin-shell bounds via parallel coupling},
arXiv:2507.15495v2, 2026.

\bibitem{KlartagLehecPolylog}
B.~Klartag and J.~Lehec,
\emph{Bourgain's slicing problem and KLS isoperimetry up to polylog},
Geom. Funct. Anal. \textbf{32} (2022), no.~5, 1134-1159.

\bibitem{KlartagLehecSurvey}
B.~Klartag and J.~Lehec,
\emph{Isoperimetric inequalities in high-dimensional convex sets},
Bull. Amer. Math. Soc. \textbf{62} (2025), 575-642.

\bibitem{HartzoulakiThesis}
M.~Hartzoulaki,
\emph{Probabilistic methods in the theory of convex bodies},
Ph.D. thesis, University of Crete, March 2003.

\bibitem{LeeVempala}
Y.~T.~Lee and S.~Vempala,
\emph{Eldan's stochastic localization and the KLS conjecture: isoperimetry, concentration and mixing},
Ann. of Math. (2) \textbf{199} (2024), no.~3, 1043-1092.

\bibitem{Letwin}
B.~Letwin,
\emph{The KLS constant is $O(\log^{1/4}n)$},
arXiv:2607.24164, 2026.

\bibitem{Milman}
E.~Milman,
\emph{On the mean-width of isotropic convex bodies and their associated $L_p$-centroid bodies},
Int. Math. Res. Not. IMRN (2015), no.~11, 3408-3423.

\bibitem{PaourisConcentration}
G.~Paouris,
\emph{Concentration of mass on convex bodies},
Geom. Funct. Anal. \textbf{16} (2006), no.~5, 1021--1049.

\bibitem{PaourisPathak}
G.~Paouris and R.~Pathak,
\emph{Optimal mean width and metric entropy estimates for convex
bodies},
arXiv:2607.29522, 2026.

\bibitem{Pisier}
G.~Pisier,
\emph{The Volume of Convex Bodies and Banach Space Geometry},
Cambridge Tracts in Mathematics, Vol.~94,
Cambridge University Press, Cambridge, 1989.

\bibitem{PivovarovMeanWidth}
P.~Pivovarov,
\emph{On the volume of caps and bounding the mean-width of an
isotropic convex body},
Math. Proc. Cambridge Philos. Soc. \textbf{149} (2010), no.~2,
317--331.

\bibitem{ReisRothvoss}
V.~Reis and T.~Rothvoss,
\emph{The subspace flatness conjecture and faster integer programming},
arXiv:2303.14605v5, 2026.

\bibitem{Rudelson}
M.~Rudelson,
\emph{Distances between non-symmetric convex bodies and the $MM^*$-estimate},
Positivity \textbf{4} (2000), no.~2, 161-178.

\bibitem{VritsiouRegularEllipsoids}
B.-H.~Vritsiou,
\emph{Regular ellipsoids and a Blaschke--Santal\'o-type inequality
for projections of non-symmetric convex bodies},
J. Funct. Anal. \textbf{286} (2024), no.~11,
Paper No.~110414, 40~pp.

\end{thebibliography}

\end{document}
